% section: 特殊関数

\section{特殊関数}

ここまで,等差型,等比型,階差型,階比型など,漸化式を一定の型に
帰着して解く方法を学んできた.しかし,右辺が少し複雑になると,
同じ道具が急に働かなくなる.

まず,次の三つを比べてみよう.
\[
  a_{n+1}=a_n^2,
  \qquad
  a_{n+1}=a_n^2-1,
  \qquad
  a_{n+1}=a_n^2-2.
  \tag{2.7.1}
\]

最初の漸化式は,初項 $a_1=c$ を決めれば,同じ計算を繰り返すことで
\[
  a_n=c^{2^{n-1}}
\]
と書ける.二乗が繰り返されるため,指数が $1,2,4,8,\ldots$ と倍々に
なる.

ところが $a_{n+1}=a_n^2-1$ では,この単純な形が崩れる.例えば
積の形や指数の形を探しても,
\[
  (u+u^{-1})^2-1=u^2+u^{-2}+1
\]
のように余分な項が残り,次の段階で同じ種類の表現へ戻ってこない.
ここでいう「解けない」とは,本章で学んできた初等的な型へそのまま
帰着できない,という意味である.

一方,定数を $-2$ にすると事情が変わる.和の形を使えば,
\[
  \left(u+u^{-1}\right)^2-2=u^2+u^{-2}
  \tag{2.7.2}
\]
となるからである.実際,
\[
  a_1=\alpha+\frac{1}{\alpha},
  \qquad
  a_{n+1}=a_n^2-2
\]
に対しては,
\[
  a_n=\alpha^{2^{n-1}}+\alpha^{-2^{n-1}}
\]
が成り立つ.三つの式は見た目がよく似ているが,閉じた表現が保たれるか
どうかは,定数項のわずかな違いに左右される.

\begin{quote}
ここで扱うのは,一般の漸化式がいつでも解けるという話ではない.
一見すると解けそうにない式の中に,三角関数,対称式,接線近似などの
別の数学的対象と偶然一致するものがある.その一致を見抜くことで,
初等的な型に還元できない漸化式を解いていく.
\end{quote}

\subsection{三角関数型}

三角関数と置くことで,数列の漸化式が角度の漸化式へ変わることがある.
その仕組みを確認するために,まず加法定理を振り返ろう.
\begin{align*}
  \sin(\alpha+\beta)
    &=\sin\alpha\cos\beta+\cos\alpha\sin\beta,\\
  \cos(\alpha+\beta)
    &=\cos\alpha\cos\beta-\sin\alpha\sin\beta,\\
  \tan(\alpha+\beta)
    &=\frac{\tan\alpha+\tan\beta}
            {1-\tan\alpha\tan\beta}.
\end{align*}
ここで $\alpha=\beta=x$ とすれば,倍角公式
\begin{align*}
  \sin 2x&=2\sin x\cos x,\\
  \cos 2x&=2\cos^2x-1=1-2\sin^2x,\\
  \tan 2x&=\frac{2\tan x}{1-\tan^2x}
\end{align*}
を得る.また,
\begin{align*}
  \sin^2x&=\frac{1-\cos 2x}{2},\\
  \cos^2x&=\frac{1+\cos 2x}{2},\\
  \tan^2x&=\frac{1-\cos 2x}{1+\cos 2x}
\end{align*}
である.三倍角公式も後で使う.
\begin{align*}
  \sin 3x&=3\sin x-4\sin^3x,\\
  \cos 3x&=4\cos^3x-3\cos x.
\end{align*}

\textbf{例1: 余弦の倍角公式を使う場合}

次の漸化式を考える.
\[
  a_1=\frac{\sqrt{3}}{2},
  \qquad
  a_{n+1}=2a_n^2-1.
  \tag{2.7.3}
\]
右辺は $\cos 2x=2\cos^2x-1$ と一致している.さらに,
初項は
\[
  a_1=\cos\frac{\pi}{6}
\]
と書ける.そこで $a_n=\cos\theta_n$ と置くと,
\[
  \theta_1=\frac{\pi}{6},
  \qquad
  \theta_{n+1}=2\theta_n
\]
と考えられる.従って候補は
\[
  a_n=\cos\left(2^{n-1}\frac{\pi}{6}\right)
  \tag{2.7.4}
\]
である.

この候補を帰納法で確認する.\,$n=1$ では初項に一致する.ある $n$ で
\[
  a_n=\cos\left(2^{n-1}\frac{\pi}{6}\right)
\]
が成り立つとすると,
\begin{align*}
  a_{n+1}
    &=2a_n^2-1\\
    &=2\cos^2\left(2^{n-1}\frac{\pi}{6}\right)-1\\
    &=\cos\left(2^n\frac{\pi}{6}\right).
\end{align*}
これは $n+1$ の式である.よって,
\[
  a_n=\cos\left(2^{n-1}\frac{\pi}{6}\right)
\]
が示された.

この解法の本質は,二次式を無理に展開することではない.数列の値を
角度の関数として表すことで,二次式の計算を角度の2倍という単純な
操作に翻訳しているのである.

\textbf{例2: 正接の倍角公式を使う場合}

今度は
\[
  a_1=\tan\frac{\pi}{12}=2-\sqrt{3},
  \qquad
  a_{n+1}=\frac{2a_n}{1-a_n^2}
  \tag{2.7.5}
\]
を考える.右辺は正接の倍角公式そのものである.したがって,
\[
  a_n=\tan\left(2^{n-1}\frac{\pi}{12}\right)
  \tag{2.7.6}
\]
と予想できる.

実際,\,$n=1$ で初項に一致し,帰納法の段階では
\begin{align*}
  a_{n+1}
    &=\frac{2a_n}{1-a_n^2}\\
    &=\frac{2\tan(2^{n-1}\pi/12)}
            {1-\tan^2(2^{n-1}\pi/12)}\\
    &=\tan\left(2^n\frac{\pi}{12}\right)
\end{align*}
となる.ただし,分母が $0$ になると正接は定義できない.
三角関数型では,置換を思いつくだけでなく,途中で定義できなくなる
項がないかも確認しなければならない.

余弦の例と正接の例は,同じ方針を共有している.すなわち,
\[
  a_n=\varphi(\theta_n),
  \qquad
  \varphi(m\theta)=F_m(\varphi(\theta))
\]
となる関数 $\varphi$ と整数倍公式を探し,角度の漸化式を解いている.

ここで現れた多項式
\[
  2x^2-1,
  \qquad
  4x^3-3x
\]
は偶然の産物ではない.チェビシェフ多項式 $T_m(x)$ を
\[
  T_m(\cos\theta)=\cos(m\theta)
  \tag{2.7.7}
\]
で定義すれば,
\[
  T_0(x)=1,
  \qquad
  T_1(x)=x,
  \qquad
  T_{m+1}(x)=2xT_m(x)-T_{m-1}(x)
\]
となる.例えば $T_2(x)=2x^2-1$ である.
したがって,余弦の例は
\[
  a_{n+1}=T_2(a_n),
  \qquad
  a_n=T_{2^{n-1}}(a_1)
\]
と書き直せる.三角関数の公式から始めて多項式の族へ進む道が,
付録のチェビシェフ多項式へつながっている.

\subsection{基本対称式型}

三角関数を使わなくても,指数の組み合わせを保つ恒等式から漸化式を
解ける場合がある.その出発点は
\[
  \left(\alpha-\frac{1}{\alpha}\right)^2
  =\alpha^2+\frac{1}{\alpha^2}-2
  \tag{2.7.8}
\]
である.

この式から分かるのは,差の形と和の形が二乗によって入れ替わる,
ということである.そこで,漸化式を解くためには,次の項でも同じ種類の
表現が保たれる形を選ぶ必要がある.
\[
  b_1=\alpha+\frac{1}{\alpha},
  \qquad
  b_{n+1}=b_n^2-2.
  \tag{2.7.9}
\]
このとき,
\[
  b_n=\alpha^{2^{n-1}}+\alpha^{-2^{n-1}}
  \tag{2.7.10}
\]
である.実際,\,$u=\alpha^{2^{n-1}}$ と書けば,
\begin{align*}
  b_{n+1}
    &=\left(u+u^{-1}\right)^2-2\\
    &=u^2+u^{-2}\\
    &=\alpha^{2^n}+\alpha^{-2^n}.
\end{align*}
指数が2倍になるたびに,和の形が保たれている.

ここで,定数を $-2$ から $-1$ に変えてみよう.
\[
  a_{n+1}=a_n^2-1.
  \tag{2.7.11}
\]
同じ置換を試すと,
\begin{align*}
  a_n^2-1
    &=\left(u+u^{-1}\right)^2-1\\
    &=u^2+u^{-2}+1,
\end{align*}
となる.しかし,次の項として欲しい形は $u^2+u^{-2}$ である.
余分な $1$ が残るため,同じ対称式の族の中に戻ってこない.

\begin{quote}
漸化式の $-2$ と $-1$ は,見た目には小さな違いである.
しかし,対称式を保つ恒等式にとっては,その違いが決定的になる.
「ほとんど同じ式だから,ほとんど同じ方法で解ける」とは限らない.
解法が成立するかどうかは,係数の大きさではなく,恒等式が完全に一致
するかどうかで決まる.
\end{quote}

同じ考え方は,三次式にも現れる.
\[
  \left(\alpha+\frac{1}{\alpha}\right)^3
  -3\left(\alpha+\frac{1}{\alpha}\right)
  =\alpha^3+\frac{1}{\alpha^3}.
\]
このように,対称な組み合わせを選ぶと,べき乗の複雑な展開が
再び同じ種類の式へ戻る.ただし,恒等式が少しでもずれると,その閉じた
構造は失われる.これが基本対称式型の面白さである.

\subsection{ニュートン法型}

これまでの例では,数列の一般項を求めることを目標にしてきた.
しかし,方程式の解を正確な式で表すよりも,解に近づく数列を作る方が
自然な場合もある.ここで登場するのがニュートン法である.

方程式
\[
  f(x)=x^2-2=0
\]
の正の解 $\sqrt{2}$ を求めたいとする.いまの近似値を $x=a_n$ とし,
点 $(a_n,f(a_n))$ における接線を引く.その接線の傾きは
$f'(a_n)=2a_n$ だから,接線の方程式は
\[
  y=f(a_n)+2a_n(x-a_n)
\]
である.\,$y=0$ と置いて $x$ 軸との交点を求めると,
\begin{align*}
  a_{n+1}
    &=a_n-\frac{f(a_n)}{f'(a_n)}\\
    &=a_n-\frac{a_n^2-2}{2a_n}\\
    &=\frac{a_n^2+2}{2a_n}.
\end{align*}
これが
\[
  a_{n+1}=\frac{a_n^2+2}{2a_n}
  \tag{2.7.12}
\]
の意味である.

例えば $a_1=2$ とすると,
\[
  a_2=\frac{3}{2},
  \qquad
  a_3=\frac{17}{12},
  \qquad
  a_4=\frac{577}{408}
\]
となり,これらは $\sqrt{2}$ に急速に近づく.この速さは,誤差を
$e_n=a_n-\sqrt{2}$ と置くと見える.
\begin{align*}
  e_{n+1}
    &=\frac{a_n^2+2}{2a_n}-\sqrt{2}\\
    &=\frac{(a_n-\sqrt{2})^2}{2a_n}\\
    &=\frac{e_n^2}{2a_n}.
\end{align*}
誤差がほぼ二乗されるため,正しい近似に近づくと桁数が一気に増える.

この漸化式の解き方は,一般項を閉じた式にすることではない.
「接線を使って,次の近似値をどう設計するか」を考え,その設計が
どのような速さで誤差を減らすかを調べている.付録では,この考えを
漸化式の収束,固定点,誤差評価,微分方程式の数値解法へ広げる.

\subsection{初項依存型}

最後に,漸化式の形だけでなく,初項の選び方そのものが問題を左右する
例を見よう.
\[
  a_1=c,
  \qquad
  a_{n+1}=\frac{a_n^2-1}{2}.
  \tag{2.7.13}
\]
初項を記号 $c$ のままにすると,
\begin{align*}
  a_2&=\frac{c^2-1}{2},\\
  a_3&=\frac{c^4-2c^2-3}{8},\\
  a_4&=\frac{c^8-4c^6-2c^4+12c^2-55}{128}.
\end{align*}
一般に $a_n$ は $c$ の次数 $2^{n-1}$ の多項式になる.初項を一般の
文字のまま扱おうとすると,項を一つ進めるたびに次数が倍になり,
三角関数型や基本対称式型のような固定された表現が見えてこない.

特に $c=2$ と固定すると,
\[
  a_1=2,
  \qquad
  a_2=\frac{3}{2},
  \qquad
  a_3=\frac{5}{8},
  \qquad
  a_4=-\frac{39}{128}
\]
となる.しかし,この列は直前の例のように角度を2倍する形にも,
指数の差を保つ形にも,接線近似の形にもならない.

ここでいう「$c=2$ 以外では解けない」とは,他の初項の数列が存在しない
という意味ではない.また,定数列を探すだけなら,
\[
  c=1+\sqrt{2},
  \qquad
  c=1-\sqrt{2}
\]
では $a_{n+1}=a_n=c$ となる.大切なのは,ある初項を選んだときだけ
特別な恒等式や置換が働くことがあり,初項を一般化するとその偶然の
構造が失われる,という事実である.

この例では,少なくとも本節で扱った初等的な道具からは一般項を得る
決まった方法が出てこない.それでも漸化式は数列を定めており,
\[
  a_n=f^{\,\circ(n-1)}(c),
  \qquad
  f(x)=\frac{x^2-1}{2}
\]
と書くことはできる.これは解答の終わりではなく,「一般項を求める」
という目標自体を見直す入口である.付録では,閉じた式が得られない
場合にも,軌道,極限,近似,行列,母関数など別の言葉で数列を理解する
方法を考える.
